\documentclass[10pt,a4paper]{amsart}
\usepackage{amsmath,amssymb,amsthm,graphicx}
% Retained preamble for the Stacks Project chapter snapshot in this directory.
% The chapter's own preamble.tex is a copy of the shared Stacks Project preamble; the block
% below is sliced MECHANICALLY, line for line, from it. Nothing was added, removed or
% redefined; the wrapper class, the xr-hyper/hyperref wiring and the externaldocument list
% are omitted because the assembler supplies the document class.
%
% lines taken from preamble.tex: 3, 5, 9, 11, 12, 156, 157, 158, ... (51 lines total)
%
% package coverage restored from the chapter's own preamble (2026-10-13): xy, xy-2cell

\usepackage[all]{xy}
\xyoption{2cell}
\UseAllTwocells

\newenvironment{reference}{\comment}{\endcomment}
\newenvironment{slogan}{\comment}{\endcomment}
\newenvironment{history}{\comment}{\endcomment}
\theoremstyle{plain}
\newtheorem{theorem}[subsection]{Theorem}
\newtheorem{proposition}[subsection]{Proposition}
\newtheorem{lemma}[subsection]{Lemma}
\theoremstyle{definition}
\newtheorem{definition}[subsection]{Definition}
\newtheorem{example}[subsection]{Example}
\newtheorem{exercise}[subsection]{Exercise}
\newtheorem{situation}[subsection]{Situation}
\theoremstyle{remark}
\newtheorem{remark}[subsection]{Remark}
\newtheorem{remarks}[subsection]{Remarks}
\numberwithin{equation}{subsection}
\def\lim{\mathop{\mathrm{lim}}\nolimits}
\def\colim{\mathop{\mathrm{colim}}\nolimits}
\def\Spec{\mathop{\mathrm{Spec}}}
\def\Hom{\mathop{\mathrm{Hom}}\nolimits}
\def\Ext{\mathop{\mathrm{Ext}}\nolimits}
\def\SheafHom{\mathop{\mathcal{H}\!\mathit{om}}\nolimits}
\def\SheafExt{\mathop{\mathcal{E}\!\mathit{xt}}\nolimits}
\def\Sch{\mathit{Sch}}
\def\Mor{\mathop{\mathrm{Mor}}\nolimits}
\def\Ob{\mathop{\mathrm{Ob}}\nolimits}
\def\Sh{\mathop{\mathit{Sh}}\nolimits}
\def\NL{\mathop{N\!L}\nolimits}
\def\CH{\mathop{\mathrm{CH}}\nolimits}
\def\proetale{{pro\text{-}\acute{e}tale}}
\def\etale{{\acute{e}tale}}
\def\QCoh{\mathit{QCoh}}
\def\Ker{\mathop{\mathrm{Ker}}}
\def\Im{\mathop{\mathrm{Im}}}
\def\Coker{\mathop{\mathrm{Coker}}}
\def\Coim{\mathop{\mathrm{Coim}}}
\DeclareMathSymbol{\boxtimes}{\mathbin}{AMSa}{"02}
\def\QCohstack{\mathcal{QC}\!\mathit{oh}}
\def\Cohstack{\mathcal{C}\!\mathit{oh}}
\def\Spacesstack{\mathcal{S}\!\mathit{paces}}
\def\Quotfunctor{\mathrm{Quot}}
\def\Hilbfunctor{\mathrm{Hilb}}
\def\Curvesstack{\mathcal{C}\!\mathit{urves}}
\def\Polarizedstack{\mathcal{P}\!\mathit{olarized}}
\def\Complexesstack{\mathcal{C}\!\mathit{omplexes}}
\def\Pic{\mathop{\mathrm{Pic}}\nolimits}
\def\Picardstack{\mathcal{P}\!\mathit{ic}}
\def\Picardfunctor{\mathrm{Pic}}
\def\Deformationcategory{\mathcal{D}\!\mathit{ef}}
\usepackage{cleveref}
\crefname{originallemma}{original lemma}{original lemmas}
\Crefname{originallemma}{Original lemma}{Original lemmas}
\makeatletter
\newlabel{categories-lemma-functorial-limit}{{002L}{source}{}{sourcepointer-1}{}}
\newlabel{categories-lemma-functorial-limit@cref}{{[originallemma][0][]002L}{[1][0][]source}{}{}{}}
\makeatother
\crefname{originallemma}{original lemma}{original lemmas}
\Crefname{originallemma}{Original lemma}{Original lemmas}
\makeatletter
\newlabel{sites-lemma-adjoints-pu}{{00XH}{source}{}{sourcepointer-2}{}}
\newlabel{sites-lemma-adjoints-pu@cref}{{[originallemma][0][]00XH}{[1][0][]source}{}{}{}}
\makeatother
\crefname{originallemma}{original lemma}{original lemmas}
\Crefname{originallemma}{Original lemma}{Original lemmas}
\makeatletter
\newlabel{sites-section-more-functoriality-PSh}{{00XF}{source}{}{sourcepointer-3}{}}
\newlabel{sites-section-more-functoriality-PSh@cref}{{[originallemma][0][]00XF}{[1][0][]source}{}{}{}}
\makeatother
\title{Stacks, Sheaves on Algebraic Stacks: lemma-2-morphisms-presheaves}
\author{The Stacks Project Authors}
\date{stacks-project:stacks-sheaves.tex:local-snapshot-20261012; GNU FDL 1.2 or later}
\begin{document}
\maketitle
\noindent\textit{Editorial note.} Statement span is lines 439-454 of sources/stacks-algebraic-stacks-sheaves/pinned-original-source.tex; proof span is lines 456-523, the complete author proof (66 body lines).\par\medskip
\noindent\textit{Editorial note.} Navigation stubs: 3 cited-but-not-included labels are supplied as non-load-bearing pointers with their pinned tags where the snapshot tag index has them.\par\medskip
\noindent\textit{Editorial note.} extra\_packages is empty on purpose: the retained preamble already loads xy with options and a second declaration clashes in the assembler wrapper.\par\medskip
\noindent\textit{Editorial note.} Text is retained verbatim, including the author's own typos and placeholders.\par\medskip
\noindent\textit{Editorial note.} Queue index 7 of queue\_algebraic-stacks-sheaves.json (pb=2729, stmt 439-454, proof 456-523).\par\medskip
\noindent\textit{Allowed standard prerequisites (editorial).}\begin{itemize}
\item Explicit prerequisite boundary. The cited inputs are prerequisites; the argument in the retained span is not replaced by any cited result.
\end{itemize}\medskip
\ifdefined\hypertarget\hypertarget{sourcepointer-1}{}\fi
\noindent\textit{Source reference.} Original lemma 002L: sources/stacks-algebraic-stacks-sheaves/pinned-original-source.tex, source lines 523--523. Cited by the retained statement or proof but not part of the collected spans; supplied only as a navigation pointer. The argument does not depend on its content.\par\medskip
\ifdefined\hypertarget\hypertarget{sourcepointer-2}{}\fi
\noindent\textit{Source reference.} Original lemma 00XH: sources/stacks-algebraic-stacks-sheaves/pinned-original-source.tex, source lines 523--523. Cited by the retained statement or proof but not part of the collected spans; supplied only as a navigation pointer. The argument does not depend on its content.\par\medskip
\ifdefined\hypertarget\hypertarget{sourcepointer-3}{}\fi
\noindent\textit{Source reference.} Original lemma 00XF: sources/stacks-algebraic-stacks-sheaves/pinned-original-source.tex, source lines 523--523. Cited by the retained statement or proof but not part of the collected spans; supplied only as a navigation pointer. The argument does not depend on its content.\par\medskip
\begin{lemma}
\label{lemma-2-morphisms-presheaves}
Let $f, g : \mathcal{X} \to \mathcal{Y}$ be $1$-morphisms of categories
fibred in groupoids over $(\Sch/S)_{fppf}$. Let $t : f \to g$
be a $2$-morphism of categories fibred in groupoids over
$(\Sch/S)_{fppf}$. Assigned to $t$ there are canonical
isomorphisms of functors
$$
t^p : g^p \longrightarrow f^p
\quad\text{and}\quad
{}_pt : {}_pf \longrightarrow {}_pg
$$
which compatible with adjointness of $(f^p, {}_pf)$ and
$(g^p, {}_pg)$ and with
vertical and horizontal composition of $2$-morphisms.
\end{lemma}

\begin{proof}
Let $\mathcal{G}$ be a presheaf on $\mathcal{Y}$. Then
$t^p : g^p\mathcal{G} \to f^p\mathcal{G}$ is given by the family
of maps
$$
g^p\mathcal{G}(x) = \mathcal{G}(g(x))
\xrightarrow{\mathcal{G}(t_x)}
\mathcal{G}(f(x)) = f^p\mathcal{G}(x)
$$
parametrized by $x \in \Ob(\mathcal{X})$. This makes sense as
$t_x : f(x) \to g(x)$ and $\mathcal{G}$ is a contravariant functor.
We omit the verification that this is compatible with restriction
mappings.

\medskip\noindent
To define the transformation ${}_pt$ for $y \in \Ob(\mathcal{Y})$
define ${}_y^f\mathcal{I}$, resp.\ ${}_y^g\mathcal{I}$ to be the category
of pairs $(x, \psi : f(x) \to y)$, resp.\ $(x, \psi : g(x) \to y)$, see
Sites, Section \ref{sites-section-more-functoriality-PSh}.
Note that $t$ defines a functor
${}_yt : {}_y^g\mathcal{I} \to {}_y^f\mathcal{I}$
given by the rule
$$
(x, g(x) \to y) \longmapsto (x, f(x) \xrightarrow{t_x} g(x) \to y).
$$
Note that for $\mathcal{F}$ a presheaf on $\mathcal{X}$ the composition
of ${}_yt$ with $\mathcal{F} : {}_y^f\mathcal{I}^{opp} \to \textit{Sets}$,
$(x, f(x) \to y) \mapsto \mathcal{F}(x)$ is equal to
$\mathcal{F} : {}_y^g\mathcal{I}^{opp} \to \textit{Sets}$. Hence by
Categories, Lemma \ref{categories-lemma-functorial-limit}
we get for every $y \in \Ob(\mathcal{Y})$ a canonical map
$$
({}_pf\mathcal{F})(y) = \lim_{{}_y^f\mathcal{I}} \mathcal{F}
\longrightarrow
\lim_{{}_y^g\mathcal{I}} \mathcal{F} = ({}_pg\mathcal{F})(y)
$$
We omit the verification that this is compatible with restriction
mappings. An alternative to this direct construction is to define
${}_pt$ as the unique map compatible with the adjointness properties
of the pairs $(f^p, {}_pf)$ and $(g^p, {}_pg)$ (see below). This also
has the advantage that one does not need to prove the compatibility.

\medskip\noindent
Compatibility with adjointness of $(f^p, {}_pf)$ and $(g^p, {}_pg)$ means
that given presheaves $\mathcal{G}$ and $\mathcal{F}$ as above we have
a commutative diagram
$$
\xymatrix{
\Mor_{\textit{PSh}(\mathcal{X})}(f^p\mathcal{G}, \mathcal{F})
\ar@{=}[r] \ar[d]_{- \circ t^p} &
\Mor_{\textit{PSh}(\mathcal{Y})}(\mathcal{G}, {}_pf\mathcal{F})
\ar[d]^{{}_pt \circ -} \\
\Mor_{\textit{PSh}(\mathcal{X})}(g^p\mathcal{G}, \mathcal{F})
\ar@{=}[r] &
\Mor_{\textit{PSh}(\mathcal{Y})}(\mathcal{G}, {}_pg\mathcal{F})
}
$$
Proof omitted. Hint: Work through the proof of
Sites, Lemma \ref{sites-lemma-adjoints-pu}
and observe the compatibility from the explicit description of the
horizontal and vertical maps in the diagram.

\medskip\noindent
We omit the verification that this is compatible with vertical and horizontal
compositions. Hint: The proof of this for $t^p$ is straightforward and
one can conclude that this holds for the ${}_pt$ maps using compatibility
with adjointness.
\end{proof}

\end{document}
